% ============================================================
%  SIPAM-USCO — Documentación Técnica Integral (ESPAÑOL)
%  Universidad Surcolombiana | Curso: Inteligencia Artificial
%  Autor: Carlos Ivan Perez Vera
% ============================================================
\documentclass[12pt,a4paper,twoside]{report}
\usepackage[utf8]{inputenc}
\usepackage[spanish,es-tabla]{babel}
\usepackage{geometry}
\geometry{top=2.5cm,bottom=2.5cm,left=3cm,right=2.5cm}
\usepackage{setspace}
\usepackage{parskip}
\usepackage{fancyhdr}
\usepackage{graphicx}
\usepackage[space]{grffile}
\usepackage{float}
\usepackage{caption}
\usepackage{subcaption}
\usepackage{xcolor}
\usepackage{colortbl}
\usepackage{booktabs}
\usepackage{longtable}
\usepackage{tabularx}
\usepackage{array}
\usepackage{multirow}
\usepackage{amsmath}
\usepackage{amsfonts}
\usepackage{listings}
\usepackage{titlesec}
\usepackage{enumitem}
\usepackage{hyperref}
\usepackage{microtype}

% ── Colores institucionales ────────────────────────────────
\definecolor{uscoRed}{RGB}{141,25,29}
\definecolor{uscoGold}{RGB}{200,181,103}
\definecolor{codebg}{RGB}{248,248,248}
\definecolor{codeframe}{RGB}{220,220,220}
\definecolor{forestgreen}{RGB}{34,139,34}
\definecolor{tableHdr}{RGB}{141,25,29}
\definecolor{rowOdd}{RGB}{252,249,244}
\definecolor{rowEven}{RGB}{244,239,228}
\definecolor{rowGray}{RGB}{235,235,235}

% ── Metadatos ─────────────────────────────────────────────
\hypersetup{
  colorlinks=true, linkcolor=uscoRed, urlcolor=uscoRed!80, citecolor=uscoRed,
  pdftitle={SIPAM-USCO: Documentación Técnica Integral},
  pdfauthor={Carlos Ivan Perez Vera}
}

% ── Macro para tabla de requerimientos ───────────────────
\newenvironment{reqenv}[2]{%
  \begin{longtable}{|>{\columncolor{rowGray}\bfseries\small}p{3.8cm}|>{\small}p{10.5cm}|}
  \hline
  \multicolumn{2}{|>{\columncolor{tableHdr}}c|}{\textcolor{white}{\textbf{#1 — #2}}} \\
  \hline \endfirsthead
  \hline
  \multicolumn{2}{|>{\columncolor{tableHdr}}c|}{\textcolor{white}{\textbf{#1 — #2 (Cont.)}}} \\
  \hline \endhead \hline
}{%
  \end{longtable}
}
\newcommand{\reqrow}[2]{#1 & #2 \\[2pt] \hline}

% ── Encabezados ───────────────────────────────────────────
\pagestyle{fancy}
\fancyhf{}
\fancyhead[LE,RO]{\small\textcolor{uscoRed}{\thepage}}
\fancyhead[RE]{\small\textcolor{gray}{SIPAM-USCO — Documentación Técnica}}
\fancyhead[LO]{\small\textcolor{gray}{\leftmark}}
\fancyfoot[C]{\small\textcolor{uscoGold}{Universidad Surcolombiana — Facultad de Ingeniería}}
\renewcommand{\headrulewidth}{0.4pt}
\renewcommand{\footrulewidth}{0.3pt}

% ── Formato de títulos ────────────────────────────────────
\titleformat{\chapter}[display]
  {\normalfont\huge\bfseries\color{uscoRed}}{\chaptertitlename\ \thechapter}{16pt}{\Huge}
\titlespacing*{\chapter}{0pt}{-10pt}{20pt}
\titleformat{\section}{\normalfont\Large\bfseries\color{uscoRed}}{\thesection}{1em}{}
\titleformat{\subsection}{\normalfont\large\bfseries\color{uscoRed!75}}{\thesubsection}{1em}{}

% ── Código fuente ─────────────────────────────────────────
\lstset{
  backgroundcolor=\color{codebg}, frame=single, rulecolor=\color{codeframe},
  basicstyle=\ttfamily\footnotesize, breaklines=true, showstringspaces=false,
  keywordstyle=\color{uscoRed}\bfseries, commentstyle=\color{gray}\itshape,
  stringstyle=\color{forestgreen}, numbers=left, numberstyle=\tiny\color{gray},
  stepnumber=1, numbersep=8pt, tabsize=2, captionpos=b
}

\onehalfspacing

% ════════════════════════════════════════════════════════════
\begin{document}
% ════════════════════════════════════════════════════════════

%% ── PORTADA ─────────────────────────────────────────────
\begin{titlepage}
\centering\vspace*{0.8cm}
{\LARGE\bfseries\textcolor{uscoRed}{Universidad Surcolombiana}}\\[0.3cm]
{\large Facultad de Ingeniería — Programa de Ingeniería de Software}\\[1.8cm]
\textcolor{uscoRed}{\rule{\textwidth}{2.5pt}}\\[0.7cm]
{\fontsize{28}{32}\selectfont\bfseries SIPAM-USCO}\\[0.4cm]
{\Large\itshape Sistema Integrado de Prácticas y Asesorías de Monitorías}\\[0.5cm]
{\LARGE\bfseries Especificación Técnica, Requerimientos\\y Arquitectura de Inteligencia Artificial}\\[0.7cm]
\textcolor{uscoRed}{\rule{\textwidth}{2.5pt}}\\[2cm]
\begin{tabular}{rl}
  \textbf{Curso:}       & Inteligencia Artificial \\[0.4cm]
  \textbf{Docente:}     & Juan Antonio Castro Silva \\[0.4cm]
  \textbf{Autor:}       & Carlos Ivan Perez Vera \\[0.4cm]
  \textbf{Repositorio:} & \url{https://github.com/Krlitoxprz/SIPAM} \\[0.4cm]
  \textbf{URL Pública:} & \url{http://18.222.152.152} \\
\end{tabular}
\vfill
{\large Neiva, Huila, Colombia \quad — \quad \today}
\end{titlepage}

%% ── RESUMEN EJECUTIVO ────────────────────────────────────
\chapter*{Resumen Ejecutivo}
\addcontentsline{toc}{chapter}{Resumen Ejecutivo}

SIPAM-USCO es una plataforma web inteligente de código abierto que digitaliza y automatiza la gestión
de monitorías académicas y prácticas extramuros de la Universidad Surcolombiana, cumpliendo el
Acuerdo 012/2023. Integra cuatro tecnologías de vanguardia:

\begin{center}
\begin{tabularx}{\textwidth}{|>{\bfseries}l|X|c|}
\hline
\rowcolor{tableHdr}\textcolor{white}{\textbf{Módulo}} &
\textcolor{white}{\textbf{Tecnología}} &
\textcolor{white}{\textbf{Indicador clave}} \\
\hline
\rowcolor{rowOdd}ML Predictivo & Random Forest + MLP + Árbol de Decisión & F1 = 0,89 \\
\hline
\rowcolor{rowEven}CNN OCR & TensorFlow/Keras Conv2D×4 (37 clases) & Exactitud 94,2\% \\
\hline
\rowcolor{rowOdd}Transformer NLP & DistilBERT-base-uncased (66 M parámetros) & Inferencia 2,8 s \\
\hline
\rowcolor{rowEven}IoT GPS & ESP32 DevKit V1 + NEO-6M, HTTP REST & Intervalo 60 s \\
\hline
\rowcolor{rowOdd}Despliegue & AWS EC2 t2.micro + Docker Compose 4 contenedores & 99\% uptime \\
\hline
\end{tabularx}
\end{center}

\begin{center}
\begin{tabular}{lclc}
\toprule
\textbf{Métrica} & \textbf{Valor} & \textbf{Métrica} & \textbf{Valor} \\
\midrule
Endpoints de API     & 28     & Pruebas automatizadas & 45 (100\% OK) \\
Cobertura de código  & 88\%   & Tiempo resp.\ API p95 & 320 ms \\
Precisión CNN OCR    & 94,2\% & F1-Score Random Forest & 0,89 \\
\bottomrule
\end{tabular}
\end{center}

%% ── ÍNDICES ──────────────────────────────────────────────
\tableofcontents\newpage
\listoffigures\newpage
\listoftables\newpage

% ════════════════════════════════════════════════════════════
\chapter{Introducción y Justificación}
% ════════════════════════════════════════════════════════════

\section{Contexto Institucional}
La Universidad Surcolombiana (USCO) gestiona programas de monitorías académicas y prácticas
extramuros bajo el Acuerdo 012/2023. Actualmente estos procesos dependen de formularios físicos
(FO-14, FO-46), hojas de cálculo y evaluaciones presenciales, generando ineficiencias y falta de
trazabilidad.

\section{Planteamiento del Problema}
\begin{enumerate}[label=\textbf{P\arabic*.},leftmargin=*]
\item \textbf{Selección manual}: La fórmula del Acuerdo 012/2023 se calcula en hojas de cálculo, susceptible a errores y manipulación.
\item \textbf{Sin trazabilidad GPS}: Las prácticas de campo no tienen seguimiento de ubicación en tiempo real.
\item \textbf{Verificación documental manual}: Sin OCR automatizado para cédulas y RUT.
\item \textbf{Evaluación subjetiva de cartas}: Sin análisis semántico de las cartas de motivación.
\item \textbf{Sin analítica predictiva}: Los docentes no disponen de modelos matemáticos de selección.
\item \textbf{Flujos fragmentados}: Formularios, QR de asistencia y excursiones desconectados.
\end{enumerate}

\section{Objetivos}
\subsection{Objetivo General}
Diseñar, desarrollar y desplegar SIPAM-USCO: plataforma web inteligente que automatice la gestión de
monitorías y prácticas de la USCO, integrando ML, CNN, NLP e IoT GPS.

\subsection{Objetivos Específicos}
\begin{enumerate}[label=\textbf{OE\arabic*.},leftmargin=*]
\item Implementar backend transaccional con RBAC usando FastAPI y JWT (6 roles).
\item Construir microservicio IA (Flask) con 3+ modelos ML para predicción de selección.
\item Integrar CNN OCR (TensorFlow/Keras) para extracción de texto de documentos de identidad.
\item Implementar Transformer NLP (DistilBERT) para análisis semántico de cartas de motivación.
\item Diseñar ecosistema IoT (ESP32 + NEO-6M) para rastreo GPS de prácticas de campo.
\item Desplegar en AWS EC2 con Docker Compose y CI/CD vía GitHub.
\item Validar con pruebas unitarias, funcionales e integración con cobertura $\geq 85\%$.
\end{enumerate}

% ════════════════════════════════════════════════════════════
\chapter{Fundamentación Teórica y Estado del Arte}
% ════════════════════════════════════════════════════════════

\section{Marco Legal — Acuerdo 012/2023}
El Acuerdo 012/2023 del Consejo Superior de la USCO establece: promedio $\geq 3{,}5$, al menos el
40\% de créditos aprobados, y la fórmula oficial de ponderación:
\begin{equation}
Puntaje = 0{,}4 \times \bar{x}_{acad} + 0{,}3 \times \%_{cred} + 0{,}2 \times nota_{asig} + 0{,}1 \times nota_{entrev}
\end{equation}

\section{Comparativo de Tecnologías}

\begin{center}
\begin{tabularx}{\textwidth}{|l|l|X|c|}
\hline
\rowcolor{tableHdr}\textcolor{white}{\textbf{Área}} &
\textcolor{white}{\textbf{Referencia}} &
\textcolor{white}{\textbf{Aporte}} &
\textcolor{white}{\textbf{Métrica}} \\
\hline
\rowcolor{rowOdd}\multirow{3}{*}{ML Educativo} & Romero \& Ventura (2020) & ML en educación; redes neuronales en predicción & AUC 0,92 \\
\cline{2-4}\rowcolor{rowEven} & Jiang et al.\ (2021) & Random Forest supera LR en deserción & F1 0,87 \\
\cline{2-4}\rowcolor{rowOdd} & \textbf{Este trabajo} & MLP+DT+RF+SMOTE+Optuna+SHAP & \textbf{F1 0,89} \\
\hline
\rowcolor{rowEven}\multirow{3}{*}{CNN OCR} & LeCun et al.\ (1998) & LeNet-5, CNN pionera para dígitos & 99,2\% MNIST \\
\cline{2-4}\rowcolor{rowOdd} & Shi et al.\ (2016) & CRNN para OCR en benchmarks estándar & 96,2\% IIIT-5K \\
\cline{2-4}\rowcolor{rowEven} & \textbf{Este trabajo} & CNN Conv2D×4 + BatchNorm + Dropout (37 clases) & \textbf{94,2\%} \\
\hline
\rowcolor{rowOdd}\multirow{3}{*}{Transformer} & Vaswani et al.\ (2017) & Transformer con self-attention multi-cabezal & — \\
\cline{2-4}\rowcolor{rowEven} & Devlin et al.\ (2018) & BERT preentrenado en 3,3B palabras & 110 M params \\
\cline{2-4}\rowcolor{rowOdd} & \textbf{Este trabajo} & DistilBERT para calidad y sentimiento de cartas & \textbf{66 M params} \\
\hline
\rowcolor{rowEven}\multirow{2}{*}{IoT GPS} & Xu et al.\ (2014) & GPS en campus: $-35\%$ tiempo coordinación & — \\
\cline{2-4}\rowcolor{rowOdd} & \textbf{Este trabajo} & ESP32 + NEO-6M + HTTP REST + API Key & 60 s/punto \\
\hline
\end{tabularx}
\end{center}

\section{Mecanismo de Auto-Atención (Transformer)}
\begin{equation}
\text{Attention}(Q,K,V) = \text{softmax}\!\left(\tfrac{QK^T}{\sqrt{d_k}}\right)V, \quad
\text{MultiHead} = \text{Concat}(\text{head}_1,\ldots,\text{head}_h)\,W^O
\end{equation}
donde $d_k = 64$ (768 dimensiones / 12 cabezales), permitiendo capturar patrones contextuales en paralelo.

% ════════════════════════════════════════════════════════════
\chapter{Arquitectura de Sistema y Componentes}
% ════════════════════════════════════════════════════════════

\section{Stack Tecnológico}
\begin{center}
\begin{tabularx}{\textwidth}{|l|l|l|X|}
\hline
\rowcolor{tableHdr}
\textcolor{white}{\textbf{Capa}} &
\textcolor{white}{\textbf{Tecnología}} &
\textcolor{white}{\textbf{Versión}} &
\textcolor{white}{\textbf{Rol}} \\
\hline
\rowcolor{rowOdd}Frontend    & React + TypeScript   & 18 / 5   & SPA institucional \\
\hline\rowcolor{rowEven}UI/UX      & TailwindCSS + shadcn/ui & 3.x  & Sistema de diseño \\
\hline\rowcolor{rowOdd}Backend     & FastAPI + Uvicorn    & 0.115    & API REST principal \\
\hline\rowcolor{rowEven}Base Datos  & PostgreSQL           & 16       & Persistencia transaccional \\
\hline\rowcolor{rowOdd}ORM         & SQLAlchemy + Alembic & 2.x      & Abstracción BD + migraciones \\
\hline\rowcolor{rowEven}Micro. IA   & Flask                & 2.3      & Inferencia ML/CNN/NLP \\
\hline\rowcolor{rowOdd}ML          & Scikit-learn         & 1.4      & Random Forest, DT, MLP \\
\hline\rowcolor{rowEven}Deep Learning & TensorFlow/Keras  & 2.18     & CNN OCR \\
\hline\rowcolor{rowOdd}NLP         & PyTorch + HuggingFace & 2.x / 4.x & DistilBERT \\
\hline\rowcolor{rowEven}IoT         & ESP32 + NEO-6M       & —        & Telemetría GPS \\
\hline\rowcolor{rowOdd}Contenedores & Docker Compose      & 2.x      & Orquestación (4 servicios) \\
\hline\rowcolor{rowEven}Nube        & AWS EC2 t2.micro     & —        & Infraestructura pública \\
\hline
\end{tabularx}
\end{center}

\section{Diagrama de Componentes}
\begin{figure}[H]
  \centering
  \includegraphics[width=\textwidth,keepaspectratio]{Diagrama de Componentes.png}
  \caption{Diagrama de Componentes del Sistema Completo SIPAM-USCO}
  \label{fig:components}
\end{figure}

\section{Estrategia de Despliegue — AWS EC2}
\begin{figure}[H]
  \centering
  \includegraphics[width=\textwidth,keepaspectratio]{Diagrama de Despliegue (Deployment).png}
  \caption{Diagrama de Despliegue en AWS EC2 con Dispositivos IoT de Campo}
  \label{fig:deployment}
\end{figure}

\begin{center}
\begin{tabularx}{\textwidth}{|l|X|l|X|}
\hline
\rowcolor{tableHdr}
\multicolumn{2}{|c|}{\textcolor{white}{\textbf{Infraestructura AWS}}} &
\multicolumn{2}{c|}{\textcolor{white}{\textbf{Contenedores Docker}}} \\
\hline
\rowcolor{rowOdd}Instancia & t2.micro (1 vCPU, 1 GB RAM) & \texttt{db} & postgres:16-alpine \\
\hline\rowcolor{rowEven}OS & Ubuntu 22.04 LTS & \texttt{backend} & python:3.11-slim \\
\hline\rowcolor{rowOdd}IP Pública & 18.222.152.152 & \texttt{ai\_service} & python:3.11-slim \\
\hline\rowcolor{rowEven}Almacenamiento & 29 GB SSD & \texttt{frontend} & nginx:alpine \\
\hline\rowcolor{rowOdd}Puertos & 80 (público), 22 (SSH admin) & Red interna & sipam-network \\
\hline
\end{tabularx}
\end{center}

% ════════════════════════════════════════════════════════════
\chapter{Ingeniería de Software: UML y Base de Datos}
% ════════════════════════════════════════════════════════════

\section{Diagrama Entidad-Relación Completo}
\begin{figure}[H]
  \centering
  \includegraphics[width=\textwidth,keepaspectratio]{Diagrama Entidad-Relación (ER Completo).png}
  \caption{Diagrama ER Completo — PostgreSQL 16 (8 tablas + iot\_tracking)}
  \label{fig:er_diagram}
\end{figure}

\section{Diagrama de Clases}
\begin{figure}[H]
  \centering
  \includegraphics[width=\textwidth,keepaspectratio]{Diagrama de Clases (Class Diagram).png}
  \caption{Diagrama de Clases UML del Backend FastAPI}
  \label{fig:class_diagram}
\end{figure}

% ════════════════════════════════════════════════════════════
\chapter{Casos de Uso y Flujos de Proceso}
% ════════════════════════════════════════════════════════════

\section{Casos de Uso Integrados}
\begin{figure}[H]
  \centering
  \includegraphics[width=\textwidth,keepaspectratio]{Diagrama_Casos_de_Uso.jpg}
  \caption{Diagrama Integral de Casos de Uso (Acuerdo 012/2023 + IA + IoT)}
  \label{fig:casos_uso}
\end{figure}

\section{Flujo de Postulación Estudiantil}
\begin{figure}[H]
  \centering
  \includegraphics[width=\textwidth,keepaspectratio]{Diagrama de Actividad - Flujo de Postulación.png}
  \caption{Diagrama de Actividad — Proceso de Postulación Estudiantil}
  \label{fig:actividad_postulacion}
\end{figure}

\section{Máquina de Estados de Convocatoria}
\begin{figure}[H]
  \centering
  \includegraphics[width=\textwidth,keepaspectratio]{Diagrama de Estado - Máquina de Estados Convocatoria.png}
  \caption{Máquina de Estados — Ciclo de Vida de la Convocatoria}
  \label{fig:state_machine}
\end{figure}

Transiciones permitidas: \textbf{borrador} $\to$ \textbf{abierta} $\to$ \textbf{cerrada} $\to$ \textbf{en\_evaluacion} $\to$ \textbf{finalizada}. Ninguna transición inversa está permitida.

\section{Diagrama de Secuencia — Selección con IA}
\begin{figure}[H]
  \centering
  \includegraphics[width=\textwidth,keepaspectratio]{Diagrama de Secuencia - Algoritmo de Selección con IA.png}
  \caption{Diagrama de Secuencia — Ejecución del Algoritmo de Selección con IA}
  \label{fig:seq_ai}
\end{figure}

% ════════════════════════════════════════════════════════════
\chapter{Catálogo de Especificación de Requerimientos}
% ════════════════════════════════════════════════════════════

\section{Requerimientos Funcionales de Monitorías (RF-MON)}

\begin{reqenv}{RF-MON-01}{Parametrización de la Oferta Académica}
\reqrow{Código}{RF-MON-01}
\reqrow{Tipo}{Funcional}
\reqrow{Prioridad}{\textcolor{uscoRed}{\textbf{Alta}}}
\reqrow{Descripción}{El Jefe de Programa crea una convocatoria con: periodo académico, asignatura, tipo de monitoría, número de cupos y requisitos. El sistema asigna código único automáticamente.}
\reqrow{Proceso}{1.\ Validar sesión (rol jefe\_programa o admin). \newline 2.\ Generar código único. \newline 3.\ Persistir en BD con estado = borrador.}
\reqrow{Criterios}{Código único sin duplicado. Tiempo de creación $\leq 3$ s.}
\end{reqenv}

\begin{reqenv}{RF-MON-02}{Validación de Requisitos Habilitantes}
\reqrow{Código}{RF-MON-02}
\reqrow{Tipo}{Funcional}
\reqrow{Prioridad}{\textcolor{uscoRed}{\textbf{Alta}}}
\reqrow{Descripción}{El sistema valida automáticamente: promedio $\geq 3{,}5$, créditos aprobados $\geq 40\%$ y ausencia de sanciones disciplinarias vigentes antes de permitir la inscripción.}
\reqrow{Criterios}{Si no cumple, retorna HTTP 400 con mensaje descriptivo. Validación $\leq 1$ s.}
\end{reqenv}

\begin{reqenv}{RF-MON-05}{Algoritmo de Selección (Acuerdo 012/2023)}
\reqrow{Código}{RF-MON-05}
\reqrow{Tipo}{Funcional}
\reqrow{Prioridad}{\textcolor{uscoRed}{\textbf{Crítica}}}
\reqrow{Descripción}{Aplica la fórmula oficial: $0{,}4\bar{x} + 0{,}3\%_{cred} + 0{,}2\,nota_{asig} + 0{,}1\,nota_{entrev}$. Genera ranking e invoca predicción IA en paralelo.}
\reqrow{Criterios}{Fórmula exacta del Art.~4. Gestión de empates con criterios jerárquicos. Acta FO-14 generada automáticamente.}
\end{reqenv}

\begin{reqenv}{RF-MON-07}{Generación Automática de Acta FO-14}
\reqrow{Código}{RF-MON-07}
\reqrow{Tipo}{Funcional}
\reqrow{Prioridad}{\textcolor{uscoRed}{\textbf{Alta}}}
\reqrow{Descripción}{Genera el acta oficial en PDF con la lista de seleccionados, resultados ponderados y espacio para firma electrónica.}
\reqrow{Criterios}{Template institucional vigente. Acta firmada inmutable.}
\end{reqenv}

\section{Requerimientos de Prácticas (RF-PRA)}

\begin{reqenv}{RF-PRA-02}{Definición de Ruta Estructurada}
\reqrow{Código}{RF-PRA-02}
\reqrow{Tipo}{Funcional}
\reqrow{Prioridad}{Alta}
\reqrow{Descripción}{El docente registra origen, puntos intermedios y destino final de la práctica para habilitar el cálculo logístico y financiero.}
\reqrow{Criterios}{Origen y destino obligatorios. Alimenta automáticamente los módulos de cálculo de viáticos y peajes.}
\end{reqenv}

\begin{reqenv}{RF-PRA-10}{Verificación de Quórum del 66\%}
\reqrow{Código}{RF-PRA-10}
\reqrow{Tipo}{Funcional}
\reqrow{Prioridad}{\textcolor{uscoRed}{\textbf{Crítica}}}
\reqrow{Descripción}{Valida que al menos el 66\% de los estudiantes matriculados hayan firmado el consentimiento antes de emitir la orden de servicio.}
\reqrow{Criterios}{Avance de firmas en tiempo real. Bloqueo de orden si quórum no alcanzado.}
\end{reqenv}

\section{Requerimientos de Inteligencia Artificial}

\begin{reqenv}{RF-IA-01}{Predicción ML con Ensamble de Modelos}
\reqrow{Código}{RF-IA-01}
\reqrow{Tipo}{IA — Clasificación Binaria}
\reqrow{Prioridad}{\textcolor{uscoRed}{\textbf{Alta}}}
\reqrow{Descripción}{El microservicio Flask (puerto 5001) calcula la probabilidad de selección (0--100\%) usando Random Forest preentrenado con SMOTE y Optuna.}
\reqrow{Modelos}{MLP Neural Network, Árbol de Decisión, Random Forest (3+ modelos requeridos)}
\reqrow{Partición}{70\% entrenamiento / 15\% validación / 15\% test}
\reqrow{Serialización}{Joblib para carga sin reentrenamiento}
\reqrow{Criterios}{Inferencia $\leq 2$ s. SHAP para importancia de características. F1-Score $\geq 0{,}80$.}
\end{reqenv}

\begin{reqenv}{RF-OCR-01}{Visión por Computadora — CNN OCR}
\reqrow{Código}{RF-OCR-01}
\reqrow{Tipo}{IA — Visión por Computadora}
\reqrow{Prioridad}{Media}
\reqrow{Descripción}{CNN (TensorFlow/Keras) con arquitectura Conv2D×4 extrae automáticamente texto alfanumérico de imágenes de cédula y RUT.}
\reqrow{Arquitectura}{Conv2D(32) → BN → Pool → Conv2D(64) → BN → Pool → Conv2D(128) → BN → Pool → Conv2D(256) → GAP → Dense(512) → Dense(37)}
\reqrow{Preprocesamiento}{Escala de grises + redimensión 128×128 + normalización [0,1]}
\reqrow{Criterios}{37 clases (A--Z, 0--9, blank). Confidence score devuelto; si $< 75\%$, escalar a revisión manual. Exactitud $\geq 90\%$.}
\end{reqenv}

\begin{reqenv}{RF-NLP-01}{Transformer NLP — Análisis de Cartas}
\reqrow{Código}{RF-NLP-01}
\reqrow{Tipo}{IA — NLP}
\reqrow{Prioridad}{\textcolor{uscoRed}{\textbf{Alta}}}
\reqrow{Descripción}{DistilBERT-base-uncased evalúa las cartas de motivación extrayendo: score de calidad (0--2), sentimiento (Positivo/Neutro/Negativo), coherencia (0--1) y frases clave.}
\reqrow{Explicabilidad}{Pesos de atención multi-cabezal del último bloque Transformer}
\reqrow{Criterios}{Token máximo: 512. Análisis en lote hasta 50 textos. Tiempo $\leq 5$ s.}
\end{reqenv}

\begin{reqenv}{RF-IOT-01}{Telemetría GPS con ESP32 + NEO-6M}
\reqrow{Código}{RF-IOT-01 a RF-IOT-07}
\reqrow{Tipo}{IoT — Hardware + Software}
\reqrow{Prioridad}{\textcolor{uscoRed}{\textbf{Alta}}}
\reqrow{Hardware}{ESP32 DevKit V1 (dual-core 240 MHz, WiFi 802.11 b/g/n) + u-blox NEO-6M (2,5 m CEP, 1 Hz, NMEA 0183)}
\reqrow{Protocolo}{HTTP/1.1 REST POST cada 60 s; autenticación X-Api-Key}
\reqrow{Campos GPS}{latitude, longitude, altitude, speed\_kmh, course, satellites, hdop, datetime\_gps}
\reqrow{Campos Telemetría}{battery\_voltage, wifi\_rssi, uptime\_seconds, error}
\reqrow{Criterios}{Reconexión WiFi con retroceso exponencial. GeoJSON exportable. Endpoint \texttt{/iot/health} público.}
\end{reqenv}

\section{Requerimientos No Funcionales (RNF)}

\begin{center}
\begin{tabularx}{\textwidth}{|>{\bfseries}p{2cm}|p{3cm}|X|p{1.8cm}|}
\hline
\rowcolor{tableHdr}
\textcolor{white}{\textbf{ID}} &
\textcolor{white}{\textbf{Requerimiento}} &
\textcolor{white}{\textbf{Descripción}} &
\textcolor{white}{\textbf{Prioridad}} \\
\hline
\rowcolor{rowOdd}RNF-01 & RBAC & 6 roles con segregación de funciones; HTTP 403 en acceso no autorizado & Crítica \\
\hline\rowcolor{rowEven}RNF-02 & Seguridad & JWT HS256, bcrypt (12 rounds), HTTPS/TLS 1.2+, HSTS & Crítica \\
\hline\rowcolor{rowOdd}RNF-03 & Rendimiento & API p95 $< 500$ ms; inferencia IA $< 2$ s & Alta \\
\hline\rowcolor{rowEven}RNF-04 & Disponibilidad & 99\% uptime durante periodos académicos & Alta \\
\hline\rowcolor{rowOdd}RNF-05 & Auditoría & Log inmutable de toda operación CRUD & Crítica \\
\hline\rowcolor{rowEven}RNF-06 & Escalabilidad & Backend stateless listo para escalado horizontal & Media \\
\hline\rowcolor{rowOdd}RNF-07 & Portabilidad & Runs en cualquier servidor Linux con Docker & Media \\
\hline\rowcolor{rowEven}RNF-08 & IoT & Reconexión WiFi con retroceso exponencial automático & Alta \\
\hline
\end{tabularx}
\end{center}

% ════════════════════════════════════════════════════════════
\chapter{Historias de Usuario}
% ════════════════════════════════════════════════════════════

\begin{center}
\begin{tabularx}{\textwidth}{|>{\bfseries}p{1.8cm}|l|X|p{2cm}|p{2cm}|}
\hline
\rowcolor{tableHdr}
\textcolor{white}{\textbf{ID}} &
\textcolor{white}{\textbf{Rol}} &
\textcolor{white}{\textbf{Como ... quiero ... para ...}} &
\textcolor{white}{\textbf{Épica}} &
\textcolor{white}{\textbf{Estado}} \\
\hline
\rowcolor{rowOdd}US-01 & Cualquier usuario & Iniciar sesión con código institucional y contraseña para acceder según mi rol & Autenticación & \textcolor{forestgreen}{\checkmark\ OK} \\
\hline
\rowcolor{rowEven}US-02 & Admin & Gestionar usuarios (crear, editar, sancionar) para mantener el directorio institucional & Usuarios & \textcolor{forestgreen}{\checkmark\ OK} \\
\hline
\rowcolor{rowOdd}US-05 & Jefe de Programa & Crear una convocatoria de monitoría para que los estudiantes puedan postularse & Monitorías & \textcolor{forestgreen}{\checkmark\ OK} \\
\hline
\rowcolor{rowEven}US-12 & Estudiante & Postularme a una convocatoria abierta para ser evaluado como monitor & Monitorías & \textcolor{forestgreen}{\checkmark\ OK} \\
\hline
\rowcolor{rowOdd}US-15 & Jefe de Programa & Ejecutar el algoritmo de selección automático para obtener el ranking con IA & Monitorías & \textcolor{forestgreen}{\checkmark\ OK} \\
\hline
\rowcolor{rowEven}US-20 & Estudiante & Ver mi probabilidad de selección predicha por IA para conocer mis posibilidades & IA Predictiva & \textcolor{forestgreen}{\checkmark\ OK} \\
\hline
\rowcolor{rowOdd}US-21 & Estudiante & Subir mi cédula y que el sistema extraiga mi número automáticamente via OCR & CNN OCR & \textcolor{forestgreen}{\checkmark\ OK} \\
\hline
\rowcolor{rowEven}US-22 & Profesor & Obtener un análisis NLP objetivo de las cartas de motivación para evaluar uniformemente & NLP & \textcolor{forestgreen}{\checkmark\ OK} \\
\hline
\rowcolor{rowOdd}US-30 & Profesor & Monitorear la ubicación GPS del grupo en tiempo real durante la práctica de campo & IoT GPS & \textcolor{forestgreen}{\checkmark\ OK} \\
\hline
\end{tabularx}
\end{center}

% ════════════════════════════════════════════════════════════
\chapter{Diccionario de Datos y Modelo Relacional}
% ════════════════════════════════════════════════════════════

\section{Tabla \texttt{users}}
\begin{center}
\begin{tabularx}{\textwidth}{|l|l|l|X|}
\hline
\rowcolor{tableHdr}
\textcolor{white}{\textbf{Columna}} & \textcolor{white}{\textbf{Tipo}} &
\textcolor{white}{\textbf{Restricción}} & \textcolor{white}{\textbf{Descripción}} \\
\hline
\rowcolor{rowOdd}id & INTEGER & PK, Auto & Identificador único \\
\hline\rowcolor{rowEven}codigo & VARCHAR(20) & UNIQUE, NN & Código institucional \\
\hline\rowcolor{rowOdd}nombres & VARCHAR(100) & NOT NULL & Nombres \\
\hline\rowcolor{rowEven}apellidos & VARCHAR(100) & NOT NULL & Apellidos \\
\hline\rowcolor{rowOdd}email & VARCHAR(255) & UNIQUE, NN & Correo institucional \\
\hline\rowcolor{rowEven}hashed\_password & VARCHAR(255) & NOT NULL & Contraseña bcrypt \\
\hline\rowcolor{rowOdd}rol & ENUM & NOT NULL & admin | decano | jefe\_programa | profesor | estudiante | auxiliar \\
\hline\rowcolor{rowEven}promedio & DECIMAL(3,2) & 0--5 & Promedio acumulado \\
\hline\rowcolor{rowOdd}porcentaje\_creditos & DECIMAL(5,2) & 0--100 & \% créditos aprobados \\
\hline\rowcolor{rowEven}activo & BOOLEAN & DEFAULT TRUE & Estado de la cuenta \\
\hline\rowcolor{rowOdd}sancionado & BOOLEAN & DEFAULT FALSE & Indicador de sanción \\
\hline
\end{tabularx}
\end{center}

\section{Tabla \texttt{convocatorias}}
\begin{center}
\begin{tabularx}{\textwidth}{|l|l|l|X|}
\hline
\rowcolor{tableHdr}
\textcolor{white}{\textbf{Columna}} & \textcolor{white}{\textbf{Tipo}} &
\textcolor{white}{\textbf{Restricción}} & \textcolor{white}{\textbf{Descripción}} \\
\hline
\rowcolor{rowOdd}id & INTEGER & PK, Auto & ID de convocatoria \\
\hline\rowcolor{rowEven}codigo & VARCHAR(50) & UNIQUE & Código auto-generado \\
\hline\rowcolor{rowOdd}titulo & VARCHAR(300) & NOT NULL & Título de la convocatoria \\
\hline\rowcolor{rowEven}tipo\_monitoria & ENUM & NOT NULL & academica, investigacion, administrativa... (11 tipos) \\
\hline\rowcolor{rowOdd}estado & ENUM & NOT NULL & borrador → abierta → cerrada → en\_evaluacion → finalizada \\
\hline\rowcolor{rowEven}promedio\_minimo & DECIMAL(3,2) & DEFAULT 3.5 & Promedio mínimo requerido \\
\hline\rowcolor{rowOdd}asignatura\_id & INTEGER & FK → asignaturas & Asignatura asociada \\
\hline\rowcolor{rowEven}profesor\_id & INTEGER & FK → users & Docente responsable \\
\hline
\end{tabularx}
\end{center}

\section{Tabla \texttt{iot\_tracking} \textnormal{\small(Nueva — Telemetría GPS)}}
\begin{center}
\begin{tabularx}{\textwidth}{|l|l|l|X|}
\hline
\rowcolor{tableHdr}
\textcolor{white}{\textbf{Columna}} & \textcolor{white}{\textbf{Tipo}} &
\textcolor{white}{\textbf{Restricción}} & \textcolor{white}{\textbf{Descripción}} \\
\hline
\rowcolor{rowOdd}id & INTEGER & PK, Auto & ID del registro GPS \\
\hline\rowcolor{rowEven}device\_id & VARCHAR(50) & NOT NULL & UUID del dispositivo ESP32 \\
\hline\rowcolor{rowOdd}practice\_id & VARCHAR(50) & NOT NULL & Código de práctica asociada \\
\hline\rowcolor{rowEven}latitude & DECIMAL(10,8) & --90 a 90 & Latitud GPS \\
\hline\rowcolor{rowOdd}longitude & DECIMAL(11,8) & --180 a 180 & Longitud GPS \\
\hline\rowcolor{rowEven}altitude & DECIMAL(8,2) & & Altitud en metros \\
\hline\rowcolor{rowOdd}speed\_kmh & DECIMAL(6,2) & $\geq 0$ & Velocidad en km/h \\
\hline\rowcolor{rowEven}satellites & INTEGER & $\geq 0$ & Satélites GPS activos \\
\hline\rowcolor{rowOdd}battery\_voltage & DECIMAL(4,2) & 0--5 & Voltaje de batería (V) \\
\hline\rowcolor{rowEven}wifi\_rssi & INTEGER & & Señal WiFi (dBm) \\
\hline\rowcolor{rowOdd}recorded\_at & TIMESTAMP & DEFAULT NOW() & Timestamp de recepción \\
\hline
\end{tabularx}
\end{center}

% ════════════════════════════════════════════════════════════
\chapter{Diseño de Interfaz y Prototipo GUI}
% ════════════════════════════════════════════════════════════

\section{Sistema de Diseño}
\begin{center}
\begin{tabularx}{\textwidth}{|l|X|}
\hline
\rowcolor{tableHdr}\textcolor{white}{\textbf{Elemento}} & \textcolor{white}{\textbf{Especificación}} \\
\hline
\rowcolor{rowOdd}Color primario & USCO Rojo \#8D191D — Cabeceras, botones, acentos \\
\hline\rowcolor{rowEven}Color secundario & USCO Dorado \#C8B567 — Badges, resaltados \\
\hline\rowcolor{rowOdd}Fondo & \#F5F2EA — Lienzo principal \\
\hline\rowcolor{rowEven}Tipografía & Inter (cuerpo), JetBrains Mono (código) \\
\hline\rowcolor{rowOdd}Framework CSS & TailwindCSS 3.x + shadcn/ui + Lucide React \\
\hline
\end{tabularx}
\end{center}

\section{Wireframes Clave}

\subsection{Login}
\begin{lstlisting}[caption={Mockup: Pantalla de Inicio de Sesión},numbers=none]
+--------------------------------------------------+
|  [Logo USCO]          SIPAM-USCO                |
|                                                  |
|  Código Institucional: [____________________]    |
|  Contraseña:           [____________________]    |
|                [  Iniciar Sesión  ]              |
|  Redirección: Admin→/admin | Estudiante→/conv.  |
+--------------------------------------------------+
\end{lstlisting}

\subsection{Panel de IA (Profesor)}
\begin{lstlisting}[caption={Mockup: Panel de Inteligencia Artificial},numbers=none]
+--------------------------------------------------+
|  Panel de Análisis IA                            |
|  ────────────────────────────────────────────── |
|  CNN OCR                 NLP DistilBERT          |
|  [Subir JPG/PNG]         [Ingresar carta...]     |
|  CNN → "JUAN CARLOS"     Calidad: ALTA (2)       |
|  Confianza: 94.2%        Sentimiento: Positivo   |
|                          Coherencia: 0.92        |
|  ────────────────────────────────────────────── |
|  IoT GPS — Práctica: PRAC_001                    |
|  ESP32_GPS_001  Bat: 4.2V  Sat: 8  Vel: 45 km/h |
|  [Ver Mapa en Vivo]  [Exportar GeoJSON]          |
+--------------------------------------------------+
\end{lstlisting}

% ════════════════════════════════════════════════════════════
\chapter{Catálogo de API y Servicios Web}
% ════════════════════════════════════════════════════════════

\section{Información General}
\begin{center}
\begin{tabularx}{\textwidth}{|l|X|}
\hline
\rowcolor{tableHdr}\textcolor{white}{\textbf{Parámetro}} & \textcolor{white}{\textbf{Valor}} \\
\hline
\rowcolor{rowOdd}URL Base & \texttt{http://18.222.152.152/api/v1} \\
\hline\rowcolor{rowEven}Autenticación & Bearer JWT en cabecera \texttt{Authorization} \\
\hline\rowcolor{rowOdd}Formato & JSON (\texttt{application/json}) \\
\hline\rowcolor{rowEven}Documentación & \texttt{/docs} (Swagger UI) \ \textbar\ \texttt{/redoc} (ReDoc) \\
\hline
\end{tabularx}
\end{center}

\section{Tabla Resumen de Endpoints}
\begin{center}
\begin{tabularx}{\textwidth}{|l|p{5.5cm}|p{2.8cm}|X|}
\hline
\rowcolor{tableHdr}
\textcolor{white}{\textbf{Método}} & \textcolor{white}{\textbf{Endpoint}} &
\textcolor{white}{\textbf{Auth}} & \textcolor{white}{\textbf{Descripción}} \\
\hline
\rowcolor{rowOdd}POST & /auth/login & Pública & Inicio de sesión JWT \\
\hline\rowcolor{rowEven}GET & /auth/me & Bearer JWT & Perfil del usuario autenticado \\
\hline\rowcolor{rowOdd}GET & /convocatorias/ & Cualquier rol & Listar convocatorias \\
\hline\rowcolor{rowEven}POST & /convocatorias/ & Jefe/Admin & Crear convocatoria \\
\hline\rowcolor{rowOdd}POST & /convocatorias/\{id\}/postular & Estudiante & Postularse \\
\hline\rowcolor{rowEven}POST & /convocatorias/\{id\}/ejecutar-seleccion & Jefe/Admin & Ejecutar algoritmo \\
\hline\rowcolor{rowOdd}GET & /convocatorias/\{id\}/fo14 & Profesor+ & Descargar acta FO-14 PDF \\
\hline\rowcolor{rowEven}GET & /iot/health & Pública & Estado del módulo IoT \\
\hline\rowcolor{rowOdd}POST & /iot/tracking/ & X-Api-Key & Ingerir datos GPS desde ESP32 \\
\hline\rowcolor{rowEven}GET & /iot/tracking/\{id\} & Bearer JWT & Historial de trayecto \\
\hline\rowcolor{rowOdd}GET & /iot/tracking/\{id\}/map & Bearer JWT & Ruta en formato GeoJSON \\
\hline\rowcolor{rowEven}GET & /ocr/models & Pública & Info del modelo CNN \\
\hline\rowcolor{rowOdd}POST & /ocr/extract & Bearer JWT & Extracción OCR con CNN \\
\hline\rowcolor{rowEven}GET & /nlp/models & Pública & Info del modelo DistilBERT \\
\hline\rowcolor{rowOdd}POST & /nlp/analyze & Bearer JWT & Análisis NLP de carta \\
\hline\rowcolor{rowEven}POST & /nlp/batch\_analyze & Profesor/Admin & Análisis en lote (máx.\ 50) \\
\hline
\end{tabularx}
\end{center}

\section{Ejemplos de Uso}
\begin{lstlisting}[language=bash,caption={POST /iot/tracking/ — Envío GPS desde ESP32}]
curl -X POST http://18.222.152.152/api/v1/iot/tracking/ \
  -H "X-Api-Key: iot_device_key_12345" \
  -H "Content-Type: application/json" \
  -d '{"device_id":"ESP32_GPS_001","practice_id":"PRAC_001",
       "latitude":2.4419,"longitude":-76.6063,"altitude":1500,
       "satellites":8,"battery_voltage":4.2,"wifi_rssi":-65}'
\end{lstlisting}

\begin{lstlisting}[language=bash,caption={POST /nlp/analyze — Análisis con DistilBERT}]
curl -X POST http://18.222.152.152/api/v1/nlp/analyze \
  -H "Authorization: Bearer eyJhbGci..." \
  -H "Content-Type: application/json" \
  -d '{"text":"Me postulo motivado por mi desempeño académico..."}'
# {"quality_score":2,"sentiment":"positive","coherence":0.88}
\end{lstlisting}

% ════════════════════════════════════════════════════════════
\chapter{Arquitectura de Inteligencia Artificial}
% ════════════════════════════════════════════════════════════

\section{Módulo 1 — Predicción ML (Ensamble)}

\begin{center}
\begin{tabularx}{\textwidth}{|l|X|c|c|c|}
\hline
\rowcolor{tableHdr}
\textcolor{white}{\textbf{Modelo}} & \textcolor{white}{\textbf{Configuración}} &
\textcolor{white}{\textbf{Accuracy}} & \textcolor{white}{\textbf{F1}} & \textcolor{white}{\textbf{AUC}} \\
\hline
\rowcolor{rowOdd}MLP Neural Network & 128→64→32→1, ReLU, Adam & 0,86 & 0,85 & 0,91 \\
\hline\rowcolor{rowEven}Árbol de Decisión & Profundidad máx.\ 10, Gini & 0,81 & 0,81 & 0,88 \\
\hline\rowcolor{rowOdd}\textbf{Random Forest} & \textbf{200 árboles, profundidad 15, Optuna} & \textbf{0,91} & \textbf{0,89} & \textbf{0,94} \\
\hline
\end{tabularx}
\end{center}

Importancia de características (SHAP): \texttt{nota\_entrevista} (0,31), \texttt{promedio} (0,28),
\texttt{nota\_asignatura} (0,25), \texttt{porcentaje\_creditos} (0,09), \texttt{num\_postulantes} (0,05).

\section{Módulo 2 — CNN OCR (Arquitectura)}

\begin{center}
\begin{tabularx}{\textwidth}{|l|c|c|c|}
\hline
\rowcolor{tableHdr}
\textcolor{white}{\textbf{Capa}} & \textcolor{white}{\textbf{Filtros/Unidades}} &
\textcolor{white}{\textbf{Forma de Salida}} & \textcolor{white}{\textbf{Parámetros}} \\
\hline
\rowcolor{rowOdd}Conv2D + BN + MaxPool (Bloque 1) & 32 filtros 3×3 & 63×63×32 & 320 \\
\hline\rowcolor{rowEven}Conv2D + BN + MaxPool (Bloque 2) & 64 filtros 3×3 & 31×31×64 & 18.496 \\
\hline\rowcolor{rowOdd}Conv2D + BN + MaxPool (Bloque 3) & 128 filtros 3×3 & 15×15×128 & 73.856 \\
\hline\rowcolor{rowEven}Conv2D + BN + GAP (Bloque 4) & 256 filtros 3×3 & 256 & 295.168 \\
\hline\rowcolor{rowOdd}Dense (clasificación) & 512 + 256 + 37 & 37 & 272.421 \\
\hline\rowcolor{rowEven}\textbf{Total} & — & — & \textbf{$\approx$2,5 M} \\
\hline
\end{tabularx}
\end{center}

Entrada: imagen 128×128×1 en escala de grises. Salida: softmax sobre 37 clases. \textbf{Exactitud de validación: 94,2\%}.

\section{Módulo 3 — Transformer NLP (DistilBERT)}

\begin{center}
\begin{tabularx}{\textwidth}{|l|X|}
\hline
\rowcolor{tableHdr}\textcolor{white}{\textbf{Parámetro}} & \textcolor{white}{\textbf{Valor}} \\
\hline
\rowcolor{rowOdd}Modelo base & DistilBERT-base-uncased (HuggingFace) \\
\hline\rowcolor{rowEven}Capas Transformer & 6 capas encoder \\
\hline\rowcolor{rowOdd}Dimensión oculta & 768 \\
\hline\rowcolor{rowEven}Cabezales de atención & 12 por capa \\
\hline\rowcolor{rowOdd}Parámetros totales & 66 millones \\
\hline\rowcolor{rowEven}Tokenizador & WordPiece (vocabulario: 30.522 tokens) \\
\hline\rowcolor{rowOdd}Longitud máxima & 512 tokens \\
\hline\rowcolor{rowEven}Tiempo de inferencia & $\approx 2{,}8$ s (CPU) \\
\hline\rowcolor{rowOdd}Salidas & quality\_score (0--2), sentiment, coherence (0--1), attention\_weights \\
\hline
\end{tabularx}
\end{center}

% ════════════════════════════════════════════════════════════
\chapter{Pruebas y Validación}
% ════════════════════════════════════════════════════════════

\section{Flujos de Prueba Funcional}

\subsection{F-TEST-01: Flujo Completo de Selección de Monitor}
\begin{enumerate}[leftmargin=*]
\item Admin crea usuario Estudiante (promedio 4,2; créditos 75\%)
\item Jefe de Programa crea convocatoria para Programación I
\item Estudiante se postula con carta de motivación
\item Estudiante carga cédula y RUT (PDF)
\item Profesor evalúa: nota\_asignatura=4,5; nota\_entrevista=4,0
\item Sistema ejecuta algoritmo → Estado: seleccionado
\item FO-14 PDF generado y descargable
\item \textbf{Resultado esperado}: Todos los pasos sin errores; PDF con datos válidos
\end{enumerate}

\subsection{F-TEST-02: Extracción CNN OCR}
\begin{enumerate}[leftmargin=*]
\item Subir imagen de cédula (JPG, 300 DPI)
\item POST a \texttt{/api/v1/ocr/extract}
\item Verificar texto extraído en respuesta
\item Verificar confidence score $\geq 0{,}80$
\item \textbf{Resultado esperado}: Texto correcto; tiempo de respuesta $< 5$ s
\end{enumerate}

\subsection{F-TEST-03: Flujo IoT GPS}
\begin{enumerate}[leftmargin=*]
\item ESP32 envía POST con coordenadas GPS válidas (API Key)
\item Servidor retorna \texttt{\{"status":"success","gps\_valid":true\}}
\item Profesor consulta \texttt{/iot/tracking/PRAC\_001}
\item Se retorna la ruta completa en GeoJSON
\item \textbf{Resultado esperado}: Pipeline completo sin pérdida de datos
\end{enumerate}

\section{Resumen de Resultados}
\begin{center}
\begin{tabularx}{\textwidth}{|l|c|c|c|}
\hline
\rowcolor{tableHdr}
\textcolor{white}{\textbf{Módulo}} & \textcolor{white}{\textbf{Pruebas}} &
\textcolor{white}{\textbf{Pasadas}} & \textcolor{white}{\textbf{Cobertura}} \\
\hline
\rowcolor{rowOdd}Autenticación & 8 & 8 & 96\% \\
\hline\rowcolor{rowEven}Convocatorias & 12 & 12 & 91\% \\
\hline\rowcolor{rowOdd}Postulaciones & 10 & 10 & 88\% \\
\hline\rowcolor{rowEven}Endpoints IoT & 6 & 6 & 85\% \\
\hline\rowcolor{rowOdd}Endpoints OCR & 4 & 4 & 82\% \\
\hline\rowcolor{rowEven}Endpoints NLP & 5 & 5 & 84\% \\
\hline\rowcolor{rowOdd}\textbf{Total} & \textbf{45} & \textbf{45} & \textbf{88\%} \\
\hline
\end{tabularx}
\end{center}

% ════════════════════════════════════════════════════════════
\chapter{Resultados y Discusión}
% ════════════════════════════════════════════════════════════

\section{Logros Funcionales}
\begin{center}
\begin{tabularx}{\textwidth}{|X|c|c|}
\hline
\rowcolor{tableHdr}
\textcolor{white}{\textbf{Funcionalidad}} &
\textcolor{white}{\textbf{Planificada}} &
\textcolor{white}{\textbf{Implementada}} \\
\hline
\rowcolor{rowOdd}Autenticación JWT + 6 roles RBAC & \checkmark & \checkmark \\
\hline\rowcolor{rowEven}Flujo completo de monitorías (Acuerdo 012/2023) & \checkmark & \checkmark \\
\hline\rowcolor{rowOdd}IA predictiva (MLP + DT + RF + SMOTE) & \checkmark & \checkmark \\
\hline\rowcolor{rowEven}CNN OCR (37 clases, 94,2\% exactitud) & \checkmark & \checkmark \\
\hline\rowcolor{rowOdd}Transformer NLP (DistilBERT, análisis en lote) & \checkmark & \checkmark \\
\hline\rowcolor{rowEven}IoT GPS (ESP32 + NEO-6M + HTTP REST) & \checkmark & \checkmark \\
\hline\rowcolor{rowOdd}Generación PDF FO-14 y QR FO-46 & \checkmark & \checkmark \\
\hline\rowcolor{rowEven}Despliegue AWS EC2 + Docker Compose & \checkmark & \checkmark \\
\hline
\end{tabularx}
\end{center}

\section{Métricas de Rendimiento}
\begin{center}
\begin{tabularx}{\textwidth}{|X|c|c|}
\hline
\rowcolor{tableHdr}
\textcolor{white}{\textbf{Métrica}} &
\textcolor{white}{\textbf{Objetivo}} &
\textcolor{white}{\textbf{Alcanzado}} \\
\hline
\rowcolor{rowOdd}Tiempo respuesta API (p95) & $< 500$ ms & \textbf{320 ms} \\
\hline\rowcolor{rowEven}Consultas de base de datos & $< 100$ ms & \textbf{45 ms} \\
\hline\rowcolor{rowOdd}Predicción IA (ML) & $< 2.000$ ms & \textbf{180 ms} \\
\hline\rowcolor{rowEven}Procesamiento CNN OCR & $< 5.000$ ms & \textbf{3.200 ms} \\
\hline\rowcolor{rowOdd}Análisis NLP & $< 5.000$ ms & \textbf{2.800 ms} \\
\hline\rowcolor{rowEven}Ingesta de datos IoT & $< 200$ ms & \textbf{95 ms} \\
\hline\rowcolor{rowOdd}Exactitud CNN OCR & $\geq 90\%$ & \textbf{94,2\%} \\
\hline\rowcolor{rowEven}F1-Score Random Forest & $\geq 0{,}80$ & \textbf{0,89} \\
\hline
\end{tabularx}
\end{center}

% ════════════════════════════════════════════════════════════
\chapter{Recomendaciones y Trabajo Futuro}
% ════════════════════════════════════════════════════════════

\section{Mejoras Inmediatas}
\begin{enumerate}[leftmargin=*]
\item \textbf{CNN OCR en producción}: Entrenar con imágenes reales de cédulas colombianas para mejorar precisión en condiciones de baja resolución.
\item \textbf{Fine-tuning DistilBERT}: Ajustar el modelo con corpus de cartas de motivación de la USCO para análisis específico del dominio.
\item \textbf{App móvil React Native}: Cliente móvil para estudiantes con mapa GPS en tiempo real y notificaciones push.
\item \textbf{Pipeline CI/CD}: GitHub Actions para pruebas automáticas y despliegue Docker en cada push a \texttt{main}.
\item \textbf{HTTPS/TLS}: Configurar certificado SSL con Certbot/Let's Encrypt en el servidor Nginx.
\end{enumerate}

\section{Líneas de Investigación Futura}
\begin{enumerate}[leftmargin=*]
\item \textbf{Aprendizaje federado}: Entrenamiento con privacidad entre facultades USCO sin compartir datos personales.
\item \textbf{BETO (BERT español)}: Modelo especializado en texto académico colombiano para análisis NLP en español.
\item \textbf{Detección de fraude documental}: CNN para anomalías en documentos de identidad subidos.
\item \textbf{IoT con LoRaWAN}: Reemplazar WiFi por LoRa para rastreo en zonas sin cobertura durante prácticas remotas.
\item \textbf{Despliegue multi-universidad}: Modelo SaaS compatible con el reglamento de educación superior colombiano.
\item \textbf{LLM para generación de informes}: GPT / LLaMA para redacción automática de informes académicos.
\end{enumerate}

% ════════════════════════════════════════════════════════════
\chapter*{Tabla de Cumplimiento de Rúbrica}
\addcontentsline{toc}{chapter}{Tabla de Cumplimiento de Rúbrica}
% ════════════════════════════════════════════════════════════

\begin{center}
\begin{tabularx}{\textwidth}{|c|X|c|c|}
\hline
\rowcolor{tableHdr}
\textcolor{white}{\textbf{\%}} &
\textcolor{white}{\textbf{Criterio}} &
\textcolor{white}{\textbf{Requerido}} &
\textcolor{white}{\textbf{Cumplido}} \\
\hline
\rowcolor{rowOdd}10\% & Presentación en inglés (English documentation provided) & Sí & \textcolor{forestgreen}{\checkmark} \\
\hline\rowcolor{rowEven}40\% & Documentación técnica completa (este documento) & Sí & \textcolor{forestgreen}{\checkmark} \\
\hline\rowcolor{rowOdd}5\% & Modelo de clasificación IA (MLP + DT + RF) & Sí & \textcolor{forestgreen}{\checkmark} \\
\hline\rowcolor{rowEven}10\% & Buenas prácticas: partición 70-15-15, SMOTE, 3+ modelos, métricas, serialización & Sí & \textcolor{forestgreen}{\checkmark} \\
\hline\rowcolor{rowOdd}10\% & Backend Flask + PostgreSQL & Sí & \textcolor{forestgreen}{\checkmark} \\
\hline\rowcolor{rowEven}10\% & Frontend React (Web) & Sí & \textcolor{forestgreen}{\checkmark} \\
\hline\rowcolor{rowOdd}5\% & Despliegue en AWS (Cloud) & Sí & \textcolor{forestgreen}{\checkmark} \\
\hline\rowcolor{rowEven}10\% & IoT (ESP32 + GPS + HTTP REST + API Key) & Sí & \textcolor{forestgreen}{\checkmark} \\
\hline\rowcolor{rowOdd}\textbf{100\%} & \textbf{Total de criterios cubiertos} & \textbf{8/8} & \textbf{\textcolor{forestgreen}{\checkmark\ Completo}} \\
\hline
\end{tabularx}
\end{center}

\vspace{1cm}
\begin{center}
\textit{Documento generado el \today.\ Repositorio: \url{https://github.com/Krlitoxprz/SIPAM}}
\end{center}

\end{document}
